% !TeX root = ../../main.tex

Wie im vorherigen Abschnitt schon angemerkt, sollten solche transient Failures wiederholt werden.
Dafür gibt es verschiedene Retry Policies.
Eine Retry Policy beschreibt sowohl die einzelnen Strategien, welche im Folgendem genauer erklärt werden, als auch
die Fehlererkennung, sowie die Anzahl der Retries und die Dauer des Delays~\cite{MicrosoftTransientFaults}.
PostgreSQL führt die Fehlererkennung mit internen Algorithmen durch, dessen Erläuterung im Rahmen dieser Arbeit
aufgrund ihrer Komplexität nicht möglich ist.

Es gibt verschiedene mögliche Retry Strategien, die in der Praxis Anwendung finden.
Eine ist der \emph{immediately Retry}, bei dem die Transaktion ohne zu warten direkt wiederholt wird.
Dabei wird nur in der Anzahl der möglichen Retries variiert.
Eine weitere in dieser Arbeit behandelten Strategie ist \emph{Retry after Delay}.
Bei diesem Mechanismus wird eine kurze Zeit gewartet, bis das Programm die Transaktion wiederholt.
Es ist sowohl der Delay als auch die Anzahl der Wiederholungen frei wählbar, sollte aber auf das System optimiert
sein~\cite{MicrosoftRetryPattern}.
Des Weiteren gibt es den \emph{Exponential back-off} und den \emph{Exponential back-off with jitter}.
Bei Exponential back-off handelt es sich um ein Verfahren, bei dem mit jedem erfolglosem Retry einer Transaktion der
Delay exponentiell länger wird.
Durch jitter wird der Delay zufällig zwischen einem Grundwert und dem Exponential back-off Wert gewählt.

Der Vollständigkeit halber sollten \emph{Cancel Retry} und \emph{Circuit Breaker}, als gängige praktische Methoden
genannt werden.
Wobei Cancel Retry nur eine Erweiterung für die anderen Strategien ist.
Diese Strategie testet den Fehlercode und wenn der Fehler kein transient fault ist, wird kein Retry ausgeführt.
Circuit Breaker hält wiederholte Anfragen ab, wenn diese eine hohe Chance haben fehlzuschlagen.
Dafür gibt es einen threshold\footnote{Schwelle} an fehlgeschlagenen Wiederholungen.
Wenn dieser überschritten wird, dann öffnet sich der Kreis und alle Transaktionen schlagen fehl, damit sich das
System erholen kann.
Zudem empfiehlt Microsoft, in ihrem Guide zum Umgang mit transient faults, für jede Strategie die einen Delay
verwendet, diesen zufällig zu generieren~\cite{MicrosoftTransientFaults}.
Für Retry Strategien müssen Transaktionen idempotent
\footnote{Eine Transaktion ist idempotent, wenn sie bei mehrfacher Ausführung
das gleiche Ergebnis liefert.} sein~\cite{MicrosoftRetryPattern}.
